%File: anonymous-submission-latex-2026.tex
\documentclass[letterpaper]{article} % DO NOT CHANGE THIS
\usepackage[submission]{aaai2026}  % DO NOT CHANGE THIS
\usepackage{times}  % DO NOT CHANGE THIS
\usepackage{helvet}  % DO NOT CHANGE THIS
\usepackage{courier}  % DO NOT CHANGE THIS
\usepackage[hyphens]{url}  % DO NOT CHANGE THIS
\usepackage{graphicx} % DO NOT CHANGE THIS
\urlstyle{rm} % DO NOT CHANGE THIS
\def\UrlFont{\rm}  % DO NOT CHANGE THIS
\usepackage{natbib}  % DO NOT CHANGE THIS AND DO NOT ADD ANY OPTIONS TO IT
\usepackage{caption} % DO NOT CHANGE THIS AND DO NOT ADD ANY OPTIONS TO IT
\frenchspacing  % DO NOT CHANGE THIS
\setlength{\pdfpagewidth}{8.5in} % DO NOT CHANGE THIS
\setlength{\pdfpageheight}{11in} % DO NOT CHANGE THIS
%
% These are recommended to typeset algorithms but not required. See the subsubsection on algorithms. Remove them if you don't have algorithms in your paper.
\usepackage{algorithm}
\usepackage{algorithmic}

%
% These are are recommended to typeset listings but not required. See the subsubsection on listing. Remove this block if you don't have listings in your paper.
\usepackage{newfloat}
\usepackage{listings}
\DeclareCaptionStyle{ruled}{labelfont=normalfont,labelsep=colon,strut=off} % DO NOT CHANGE THIS
\lstset{%
	basicstyle={\footnotesize\ttfamily},% footnotesize acceptable for monospace
	numbers=left,numberstyle=\footnotesize,xleftmargin=2em,% show line numbers, remove this entire line if you don't want the numbers.
	aboveskip=0pt,belowskip=0pt,%
	showstringspaces=false,tabsize=2,breaklines=true}
\floatstyle{ruled}
\newfloat{listing}{tb}{lst}{}
\floatname{listing}{Listing}
%
% Keep the \pdfinfo as shown here. There's no need
% for you to add the /Title and /Author tags.
\pdfinfo{
/TemplateVersion (2026.1)
}

% DISALLOWED PACKAGES
% \usepackage{authblk} -- This package is specifically forbidden
% \usepackage{balance} -- This package is specifically forbidden
% \usepackage{color (if used in text)
% \usepackage{CJK} -- This package is specifically forbidden
% \usepackage{float} -- This package is specifically forbidden
% \usepackage{flushend} -- This package is specifically forbidden
% \usepackage{fontenc} -- This package is specifically forbidden
% \usepackage{fullpage} -- This package is specifically forbidden
% \usepackage{geometry} -- This package is specifically forbidden
% \usepackage{grffile} -- This package is specifically forbidden
% \usepackage{hyperref} -- This package is specifically forbidden
% \usepackage{navigator} -- This package is specifically forbidden
% (or any other package that embeds links such as navigator or hyperref)
% \indentfirst} -- This package is specifically forbidden
% \layout} -- This package is specifically forbidden
% \multicol} -- This package is specifically forbidden
% \nameref} -- This package is specifically forbidden
% \usepackage{savetrees} -- This package is specifically forbidden
% \usepackage{setspace} -- This package is specifically forbidden
% \usepackage{stfloats} -- This package is specifically forbidden
% \usepackage{tabu} -- This package is specifically forbidden
% \usepackage{titlesec} -- This package is specifically forbidden
% \usepackage{tocbibind} -- This package is specifically forbidden
% \usepackage{ulem} -- This package is specifically forbidden
% \usepackage{wrapfig} -- This package is specifically forbidden
% DISALLOWED COMMANDS
% \nocopyright -- Your paper will not be published if you use this command
% \addtolength -- This command may not be used
% \balance -- This command may not be used
% \baselinestretch -- Your paper will not be published if you use this command
% \clearpage -- No page breaks of any kind may be used for the final version of your paper
% \columnsep -- This command may not be used
% \newpage -- No page breaks of any kind may be used for the final version of your paper
% \pagebreak -- No page breaks of any kind may be used for the final version of your paperr
% \pagestyle -- This command may not be used
% \tiny -- This is not an acceptable font size.
% \vspace{- -- No negative value may be used in proximity of a caption, figure, table, section, subsection, subsubsection, or reference
% \vskip{- -- No negative value may be used to alter spacing above or below a caption, figure, table, section, subsection, subsubsection, or reference
\usepackage{amsmath}
\usepackage{amssymb}
\usepackage{mathtools}
\usepackage{amsfonts}
\usepackage{booktabs}
\usepackage{multirow}
\usepackage{array}
\usepackage{placeins}
\newcommand{\algbreak}{\allowbreak} % safe break point in long math
% --- AAAI-style algorithmic cosmetics (safe) ---
\renewcommand{\algorithmicrequire}{\textbf{Input:}}
\renewcommand{\algorithmicensure}{\textbf{Output:}}
\renewcommand{\algorithmiccomment}[1]{\hfill{\footnotesize\ttfamily // #1}}
\newcommand{\algsmall}{\footnotesize}
\newcommand{\one}{\mathbf{1}}
\newcommand{\best}[1]{\textbf{#1}}
\newcommand{\second}[1]{\underline{#1}}
\makeatletter
\let\cmtflowoldthebibliography\thebibliography
\let\cmtflowoldendthebibliography\endthebibliography
\renewenvironment{thebibliography}[1]{%
  \cmtflowoldthebibliography{#1}%
  \small
  \setlength{\itemsep}{0pt}%
  \setlength{\parsep}{0pt}%
  \setlength{\parskip}{0pt}%
}{\cmtflowoldendthebibliography}
\makeatother
\setcounter{secnumdepth}{0} %May be changed to 1 or 2 if section numbers are desired.

% The file aaai2026.sty is the style file for AAAI Press
% proceedings, working notes, and technical reports.
%

% Title

% Your title must be in mixed case, not sentence case.
% That means all verbs (including short verbs like be, is, using,and go),
% nouns, adverbs, adjectives should be capitalized, including both words in hyphenated terms, while
% articles, conjunctions, and prepositions are lower case unless they
% directly follow a colon or long dash
\title{\textit{CMTFlow}: Hierarchical Portfolio Management with Controller-Guided Base Revision and Daily Refinement}
\author{
    Anonymous Authors
}
\affiliations{
    Anonymous Affiliation
}

\begin{document}

\maketitle

\begin{abstract}
Portfolio management in non-stationary markets requires more than predicting a daily portfolio weight: a strategy must decide when the current base portfolio has become stale, what replacement base should anchor the next holding segment, and how daily exposures should be refined without destroying that anchor. Existing reinforcement learning methods usually entangle these roles in a single daily allocation policy or rely on fixed rebalancing schedules, making revision either noisy or exogenous. We propose \textit{CMTFlow}, a hierarchical portfolio management framework that treats the active base portfolio as a stateful memory. An outer actor proposes a sparse candidate base, a controller learns a daily hold/switch policy by comparing the drifted current holding with this candidate, and an inner actor performs support-constrained refinement inside the selected base. The three modules are trained through fixed-segment HRL warmup, counterfactual controller policy-gradient rollouts, and controller-active joint finetuning, so the final checkpoint matches the daily decision protocol used at evaluation. Experiments on Nasdaq-100 and CSI-300 show that \textit{CMTFlow} achieves a favorable return--risk trade-off against matched traditional and deep RL baselines, obtaining the best or second-best total return, Sharpe ratio, maximum drawdown, and CR on most market-metric columns. Ablations and interpretability analyses further show that the controller is the main source of adaptive revision behavior, while the inner actor contributes local within-holding-period allocation refinement.
\end{abstract}

% Uncomment the following to link to your code, datasets, an extended version or similar.
% You must keep this block between (not within) the abstract and the main body of the paper.
% \begin{links}
%     \link{Code}{https://aaai.org/example/code}
%     \link{Datasets}{https://aaai.org/example/datasets}
%     \link{Extended version}{https://aaai.org/example/extended-version}
% \end{links}
\section{Introduction}

Portfolio management (PM) aims to allocate capital across assets under uncertainty and trading constraints to maximize long-horizon risk-adjusted returns. Classical approaches, such as mean--variance optimization \citep{MVO} and the Capital Asset Pricing Model \citep{CAMP}, laid the foundation for modern asset allocation, while heuristic rebalancing and threshold-based rebalancing rules remain widely used in practice to control portfolio deviation, turnover, and risk exposure \citep{Perold1988DynamicSF,9}. However, these approaches typically rely on static assumptions, single-period reasoning, or manually specified rebalancing schedules, and thus often adapt too slowly in non-stationary markets where return distributions shift, asset relationships evolve, and portfolio effectiveness may deteriorate unevenly over time.

Recent reinforcement learning (RL) methods provide a promising alternative by formulating PM as a sequential decision problem and directly optimizing long-horizon investment objectives. Existing studies have shown strong potential in end-to-end portfolio construction, cross-asset dependency modeling, market-aware representation learning, and interpretability \citep{4,5,6,8,alphastock}, while hierarchical or adaptive RL has also been explored to improve flexibility in portfolio selection or trading \citep{25}. Despite this progress, a key design question remains under-specified: what object should be held persistently, what object should be revised, and what object should be adjusted daily? Unified daily-allocation policies mainly output a full portfolio at every step, so medium-horizon base construction and short-horizon execution are optimized as the same action. Adaptive interval methods improve the timing of updates, but the revision target is usually a new allocation rather than an explicit active base that can be held, drifted with prices, compared with a candidate replacement, and locally refined. As a result, existing formulations often make revision timing either too noisy, because every day is a full reallocation problem, or too exogenous, because a fixed or interval-level rule determines when the portfolio is reconsidered.

These limitations point to a more fundamental challenge in PM: beyond learning better portfolio weights, an effective system must organize portfolio decisions according to their distinct roles and temporal scales. In real investment practice, a medium-horizon position often acts as a base portfolio: it is kept stable while its weights drift with prices, adjusted locally as new information arrives, and revised only when the current holding state or a concrete replacement candidate justifies a new segment. This perspective leads to a different modeling target from standard daily allocation. The policy should explicitly coordinate three coupled decisions: \emph{when} to revise the active base portfolio, \emph{what} base portfolio to hold over the next segment, and \emph{how} to perform daily refinement around that base portfolio.

We propose \textit{CMTFlow}, a hierarchical portfolio management framework built around this stateful base-portfolio view. The outer actor constructs a sparse candidate base for the next segment, the controller decides whether the drifted current base should be held or replaced by that candidate, and the inner actor performs support-constrained daily refinement inside the selected base. A key innovation is to model base revision as a trainable daily event policy rather than a fixed calendar schedule. The controller conditions on a recent market tensor, the drifted current holding, the candidate base, holding-state features, and action-comparison features such as turnover, concentration, and support overlap to output an exit probability. By separating base revision, segment-level allocation, and local refinement, \textit{CMTFlow} reframes PM as coordinated multi-scale control over a persistent portfolio memory rather than as a monolithic daily weighting problem.

Experiments on Nasdaq-100 and CSI-300 support this design. \textit{CMTFlow} achieves the highest matched total return on Nasdaq-100 and the best Sharpe ratio on CSI-300, while keeping drawdown substantially lower than high-return but volatile competitors and obtaining the second-best CR in both markets. The ablation results identify the controller as the dominant source of adaptive behavior: learned revision improves over fixed-segment HRL and fixed-window controllers, showing that the gain comes from state-dependent base revision rather than merely from adding more policy layers or rebalancing more frequently. Interpretability analyses further connect selected switch cases, fixed-window comparisons, and inner-actor tilts to the economic roles claimed by the method.

We summarize our contributions as follows:
\begin{itemize}
    \item We formulate portfolio management as coordinated control over a persistent active base portfolio, clarifying three coupled decisions that are often entangled in prior RL formulations: revision timing, segment-level base construction, and daily refinement.
    \item We propose \textit{CMTFlow}, a controller--outer--inner hierarchy in which the outer actor proposes a sparse candidate base, the controller learns whether to hold or replace the drifted active base, and the inner actor refines weights only inside the selected support.
    \item We provide matched-market experiments, component ablations, fixed-window controller comparisons, and interpretability studies on Nasdaq-100 and CSI-300, showing that state-dependent base revision improves the return--risk trade-off while the inner actor plays a complementary local-refinement role.
\end{itemize}

\section{Related Work}
\subsection{Classical Portfolio Optimization and Heuristic Rebalancing}
Modern portfolio theory is rooted in the mean--variance optimization framework \citep{MVO}, while the Capital Asset Pricing Model (CAPM) provides an equilibrium-based interpretation of asset pricing and risk compensation \citep{CAMP}. These classical approaches typically rely on distributional stationarity or single-period assumptions, which limit their applicability under non-stationary financial time series with evolving market regimes. As a result, heuristic rebalancing rules remain an important practical mechanism for portfolio adjustment.

Research on dynamic asset allocation suggests that fixed rebalancing rules can become misaligned with changing market conditions \citep{Perold1988DynamicSF}. Opportunistic rebalancing discusses condition-triggered adjustment ideas and the trade-off between deviation thresholds and transaction costs for reducing unnecessary turnover \citep{9}. Meanwhile, transaction costs play a critical role in portfolio optimization and must be modeled explicitly in the decision process \citep{10}.

\subsection{Reinforcement Learning for Portfolio Management}
Reinforcement learning (RL) provides a sequential decision-making framework for portfolio management. Early works such as PGPortfolio demonstrated the feasibility of end-to-end portfolio weight generation via policy gradient methods \citep{4}, while FinRL systematized RL-based trading and portfolio management into a reusable toolkit \citep{5}.

In terms of risk control and cost modeling, some studies incorporate maximum drawdown and other downside-risk measures into the learning objective \citep{11}, while others explore RL formulations that account for transaction costs and related portfolio constraints \citep{23}. Cost-sensitive deep RL frameworks have also been proposed to explicitly account for trading frictions in portfolio learning \citep{24}.

From the perspective of model architecture, DeepTrader incorporates market-condition representations into portfolio learning \citep{6}, while HADAPS explores hierarchical adaptive multi-asset portfolio selection \citep{25}. DeepAries is particularly relevant because it learns adaptive rebalancing interval selection for portfolio selection \citep{deeparies2025}. Regarding interpretability and cross-asset relationship modeling, AlphaStock integrates cross-asset attention with interpretability analysis to model relative-strength relationships across assets \citep{alphastock}. In terms of knowledge integration and exploration, Investor-Imitator enhances policy learning through imitation learning and knowledge extraction mechanisms \citep{8}.

The closest line of work to \textit{CMTFlow} is adaptive or hierarchical portfolio learning, including HADAPS and DeepAries. These methods improve flexibility over a single fixed allocation rule, but they usually do not explicitly maintain the active base portfolio as a persistent memory state that can be held, drifted, compared with a candidate base, or replaced by a controller. \textit{CMTFlow} differs by training a controller to decide when this base memory should be replaced by a candidate outer portfolio, while the inner actor is restricted to support-constrained tilting inside the active base. This separates revision timing, segment-level base construction, and daily refinement rather than learning only an adaptive interval or a single daily allocation policy.

\section{Problem Formulation}
We consider a long-only portfolio over \(M\) assets and \(T\) trading days. Let \(c_{i,t}\) be the close price of asset \(i\) at day \(t\), and let \(\mathbf{y}_t=(c_{1,t+1}/c_{1,t},\ldots,c_{M,t+1}/c_{M,t})^\top\) be the close-to-close price-relative vector. At each day \(t\), the investor chooses a fully invested portfolio \(\mathbf{w}_t\in\Delta^M\), where \(\Delta^M=\{\mathbf{w}\mid \mathbf{1}^\top\mathbf{w}=1,\mathbf{w}\ge0\}\).

Because prices move between two decisions, the portfolio carried into day \(t\) is the drifted previous portfolio
\[
\mathcal{D}(\mathbf{w},\mathbf{y})
=
\frac{\mathbf{y}\odot \mathbf{w}}{\mathbf{y}^\top \mathbf{w}},
\qquad
\tilde{\mathbf{w}}_t=\mathcal{D}(\mathbf{w}_{t-1},\mathbf{y}_{t-1}),
\]
where \(\odot\) denotes element-wise multiplication. If the investor rebalances from \(\tilde{\mathbf{w}}_t\) to \(\mathbf{w}_t\), the proportional transaction-cost multiplier is
\[
\mu_t
=
1-\rho\|\mathbf{w}_t-\tilde{\mathbf{w}}_t\|_1,
\]
where \(\rho\) is the transaction fee rate. The portfolio value evolves as
\[
V_{t+1} = V_t\,\mu_t\, \mathbf{y}_t^\top \mathbf{w}_t.
\]
The objective is to find a feasible decision sequence that maximizes terminal wealth:
\[
\{\mathbf{w}_t^*\}_{t=0}^{T-1}
= \arg\max_{\mathbf{w}_t\in\Delta^M} V_T .
\].

\section{Methodology}
CMTFlow formalizes multi-scale portfolio management as control over an active base-portfolio memory. A purely daily policy can react quickly but may be dominated by short-term noise; a purely low-frequency policy is stable but may fail to revise a stale position. CMTFlow therefore keeps a persistent active base portfolio, lets it drift with prices when it remains useful, replaces it when a learned controller chooses to switch, and allows daily refinement only inside the selected support.

The framework contains three coordinated components, as shown in Figure~\ref{fig:framework}. The outer actor proposes a sparse candidate base portfolio for the next segment, the controller decides whether the drifted active base should be held or replaced by that candidate, and the inner actor refines the selected base into the executed daily portfolio. The implementation uses cost-aware log returns for wealth feedback, relative-return signals for daily refinement, and counterfactual return uplift for controller learning. The controller is not a post-hoc rule attached to a separately fixed policy stack: after fixed-segment warmup and controller policy-gradient learning, the final checkpoint is obtained through controller-active joint finetuning of all three modules.

We use the following notation for this daily decision flow. The vector \(\tilde{\mathbf{w}}_t=\mathcal{D}(\mathbf{w}_{t-1},\mathbf{y}_{t-1})\) is the drifted executed holding carried into day \(t\), and \(\tilde{\mathbf{b}}_t=\mathcal{D}(\mathbf{b}_{t-1},\mathbf{y}_{t-1})\) is the drifted active base before the controller decision. The outer actor proposes a candidate base \(\mathbf{w}_t^{\mathrm{cand}}\). The controller action \(g_t\in\{0,1\}\) selects hold (\(g_t=0\)) or switch (\(g_t=1\)), producing the active base \(\mathbf{b}_t\). The inner actor then constructs a support-constrained target \(\mathbf{w}_t^{\mathrm{in}}\), and the final executed portfolio is \(\mathbf{w}_t\).

Under this active-base formulation, the investment horizon is divided by a sequence of segment boundaries
\[
0=\tau_0<\tau_1<\cdots<\tau_K\le T,
\]
where the \(k\)-th holding segment is \([\tau_k,\tau_{k+1})\). At each controller decision day \(t\), the outer actor can be invoked to produce a candidate base portfolio \(\mathbf{w}_{t}^{\mathrm{cand}}\in\Delta^M\), where \(w_{i,t}^{\mathrm{cand}}>0\) indicates that asset \(i\) is selected. Between two consecutive revision points, the accepted base portfolio is not re-optimized, but instead evolves passively with market returns according to the drift mechanism defined in the problem formulation. The controller turns the drifted base \(\tilde{\mathbf{b}}_t\) and candidate base \(\mathbf{w}_t^{\mathrm{cand}}\) into the active base portfolio
\[
\mathbf{b}_t
=
g_t\mathbf{w}_t^{\mathrm{cand}}+(1-g_t)\tilde{\mathbf{b}}_t,
\qquad
\tilde{\mathbf{b}}_t=\mathcal{D}(\mathbf{b}_{t-1},\mathbf{y}_{t-1}),
\]
where \(g_t=1\) means switching to the candidate and \(g_t=0\) means holding the drifted base. In the final evaluation, the same update is applied daily with a 30-day maximum-hold cap. After \(\mathbf{b}_t\) is selected, the inner actor produces a support-constrained refinement portfolio \(\mathbf{w}_t^{\mathrm{in}}\), and the final executed portfolio is
\[
\mathbf{w}_t=(1-\alpha)\mathbf{b}_t+\alpha\mathbf{w}_t^{\mathrm{in}},
\]
where \(\alpha\in[0,1]\) controls the strength of daily refinement.

A key property of \textit{CMTFlow} is that segment boundaries are governed by a learned event policy rather than by a fixed calendar rule alone. Since the effectiveness of the active base portfolio may deteriorate unevenly over time, base-portfolio revision should be governed by the evolving state of the held base portfolio. \textit{CMTFlow} therefore updates the segment boundary when revision is triggered by the controller, with the fixed reference horizon retained only as a maximum holding cap in the reported evaluation. Once a revision occurs, a new base portfolio is generated and a new segment begins; otherwise, the current base portfolio is carried forward and daily refinement continues within the existing segment.

\begin{figure*}[t]
\centering
\includegraphics[width=\textwidth]{figures/cmtflow_architecture_vector.pdf}
\caption{\textbf{\textit{CMTFlow} decision modules.} The controller, outer actor, and inner actor form three complementary decision modules. The controller encodes the recent market tensor and compares the drifted current holding with the outer actor's candidate base using holding-state and action-comparison features; the outer actor encodes market windows to select and allocate a sparse candidate base portfolio; and the inner actor applies support-constrained daily adjustments around the active base portfolio.}
\label{fig:framework}
\end{figure*}

\subsection{Hierarchical Allocation and Refinement}

Given the segment boundaries determined by the controller, \textit{CMTFlow} constructs portfolios through two coordinated decision layers operating at different temporal resolutions. The key idea is not merely to decompose decisions by frequency, but to separate two functionally distinct roles: \textbf{segment-level base portfolio construction} and \textbf{support-constrained daily refinement}. This separation is important because stable segment-level base allocation and short-horizon daily refinement respond to different signals and should not be optimized under the same objective at the same granularity.

\paragraph{Outer Actor: Segment-Level Base Portfolio Construction.}
At each controller decision day \(t\), the outer actor observes a state
\[
s^{\mathrm{out}}_{t}
=
\bigl\{\mathbf{X}^{\mathrm{out}}_{t},\tilde{\mathbf{w}}_{t}\bigr\},
\]
where \(\mathbf{X}^{\mathrm{out}}_{t}\) summarizes long-horizon multi-asset market information up to \(t\), and \(\tilde{\mathbf{w}}_{t}\) is the drifted executed portfolio carried into the current decision point. Based on this state, the outer actor outputs a candidate base portfolio \(\mathbf{w}^{\mathrm{cand}}_{t}\). Its role is not to react to daily fluctuations, but to construct a stable segment-level allocation that can become the active base if accepted by the controller.

Accordingly, the outer actor is formulated as a segment-level Markov decision process
\[
\mathcal{M}^{\mathrm{out}}
=
\bigl(
\mathcal{S}^{\mathrm{out}},
\mathcal{A}^{\mathrm{out}},
\mathcal{R}^{\mathrm{out}},
\gamma^{\mathrm{out}},
\pi^{\mathrm{out}}
\bigr),
\]
where \(\mathcal{S}^{\mathrm{out}}\), \(\mathcal{A}^{\mathrm{out}}\), and \(\mathcal{R}^{\mathrm{out}}\) denote the state, action, and reward spaces, \(\gamma^{\mathrm{out}}\) is the discount factor, and \(\pi^{\mathrm{out}}\) is the outer policy. In implementation, the outer actor maps the segment-level input
\[
s_{t}^{\mathrm{out}}
=
\bigl(
\mathbf{X}_{t}^{\mathrm{out}},
\tilde{\mathbf{w}}_{t}
\bigr)
\]
to an asset-wise score vector through
\[
\mathbf{q}_{t}^{\mathrm{out}}
=
\mathcal{F}_{\mathrm{out}}
\!\left(
s_{t}^{\mathrm{out}};
\theta^{\mathrm{out}}
\right)
\in\mathbb{R}^{M},
\]
where \(\mathcal{F}_{\mathrm{out}}(\cdot)\) denotes the outer scoring network, implemented by an \textbf{LSTM-HA} \citep{alphastock} encoder for long-horizon asset-wise temporal modeling, followed by a \textbf{CAAN} \citep{alphastock} module, a portfolio-state fusion layer, and an \textbf{MLP} policy head for score generation, and \(\theta^{\mathrm{out}}\) collects the corresponding learnable parameters. The output
\[
\mathbf{q}_{t}^{\mathrm{out}}
=
(q_{1,t}^{\mathrm{out}},\dots,q_{M,t}^{\mathrm{out}})^\top
\]
assigns a score to each asset in the stock universe.

To construct a sparse long-only candidate base portfolio, only the Top-\(K\) assets with the largest scores are retained. Let \(\mathcal{I}_{t}=\mathrm{TopK}(\mathbf{q}^{\mathrm{out}}_{t},K)\) denote the index set of the selected assets. The candidate base portfolio is then given by
\[
w_{i,t}^{\mathrm{cand}}
=
\begin{cases}
\dfrac{\exp(q_{i,t}^{\mathrm{out}})}
{\sum_{j\in\mathcal{I}_{t}}
\exp(q_{j,t}^{\mathrm{out}})},
& i\in\mathcal{I}_{t},\\[10pt]
0, & \text{otherwise},
\end{cases}
\]
which converts asset-wise scores into a sparse long-only segment-level candidate through score-based selection and normalized capital allocation. In this way, the outer actor transforms long-horizon market context into a candidate base portfolio rather than a day-level trading action.

The outer reward is assigned at segment boundaries from the realized cost-aware log returns accumulated over the segment:
\[
r^{\mathrm{out}}_{\tau_k}
=
\sum_{t=\tau_k}^{\tau_{k+1}-1}
\ell_t,
\qquad
\ell_t=\log\!\left(\mu_t\,\mathbf{y}_t^\top\mathbf{w}_t\right),
\]
where \(\mu_t\) is the transaction-cost multiplier defined in the problem formulation and \(\mathbf{w}_t\) is the executed portfolio after daily refinement. In the PPO buffer, these daily log returns are accumulated between two switch points and credited to the accepted outer candidate that opened the segment. Thus, the outer actor is trained from realized wealth consequences of its sparse candidate base portfolio, rather than from a handcrafted equal-weight or segment-Sharpe target. Further implementation details of the outer architecture and base-portfolio construction are provided in the supplementary material.

\paragraph{Inner Actor: Support-Constrained Daily Refinement.}
Within each active segment, the inner actor acts at every trading day \(t\) based on a short-horizon refinement state
\[
s_t^{\mathrm{in}}
=
\bigl\{\mathbf{X}^{\mathrm{in}}_t,\tilde{\mathbf{w}}_t,\mathbf{b}_t\bigr\},
\]
where \(\mathbf{X}^{\mathrm{in}}_t\) encodes recent market observations, \(\tilde{\mathbf{w}}_t\) is the current drifted executed portfolio, and \(\mathbf{b}_t\) is the active base portfolio selected by the controller. Unlike a free-form daily trading policy, the inner actor is explicitly constrained to refine only within the support set of the active base portfolio. This restriction is crucial: it preserves the outer actor's segment-level allocation role while still allowing the inner actor to exploit short-horizon return asymmetries within the current segment.

Accordingly, the inner actor is formulated as a daily Markov decision process
\[
\mathcal{M}^{\mathrm{in}}
=
\bigl(
\mathcal{S}^{\mathrm{in}},
\mathcal{A}^{\mathrm{in}},
\mathcal{R}^{\mathrm{in}},
\gamma^{\mathrm{in}},
\pi^{\mathrm{in}}
\bigr),
\]
where \(\mathcal{S}^{\mathrm{in}}\), \(\mathcal{A}^{\mathrm{in}}\), and \(\mathcal{R}^{\mathrm{in}}\) denote the state, action, and reward spaces, \(\gamma^{\mathrm{in}}\) is the discount factor, and \(\pi^{\mathrm{in}}\) is the inner policy. In implementation, the inner actor maps the day-level state \(s_t^{\mathrm{in}}\) to an asset-wise refinement score vector through
\[
\mathbf{q}_t^{\mathrm{in}}
=
\mathcal{F}_{\mathrm{in}}
\!\left(
\mathbf{X}_t^{\mathrm{in}},\tilde{\mathbf{w}}_t,\mathbf{b}_t;
\theta^{\mathrm{in}}
\right)
\in\mathbb{R}^{M},
\]
where \(\mathcal{F}_{\mathrm{in}}(\cdot)\) denotes the inner refinement network. In the implementation used for the reported experiments, this network is an asset-wise \textbf{LSTM-attention} encoder: a two-layer LSTM encodes each asset's short market window, two temporal attention layers query the sequence using the latest hidden state, and the resulting asset embedding is fused with the active base weight and drifted current weight before the actor head generates refinement scores. The output \(\mathbf{q}_t^{\mathrm{in}}=(q_{1,t}^{\mathrm{in}},\dots,q_{M,t}^{\mathrm{in}})^\top\) assigns a refinement score to each asset at day \(t\).

To preserve the structural role of the outer actor, the inner actor is constrained to refine only the support of the active base portfolio. Let
\[
m_{i,t}
=
\mathbb{I}\!\left(b_{i,t}>0\right)
\]
be the support mask induced by \(\mathbf{b}_t\). The refinement portfolio is then given by
\[
w_{i,t}^{\mathrm{in}}
=
\frac{
m_{i,t}\exp(q_{i,t}^{\mathrm{in}})
}{
\sum_{j=1}^{M} m_{j,t}\exp(q_{j,t}^{\mathrm{in}})
},
\]
which produces a support-constrained local adjustment around the active base portfolio rather than reconstructing the entire allocation from scratch. The final executed portfolio is obtained through
\[
\mathbf{w}_t
=
(1-\alpha)\mathbf{b}_t
+
\alpha\mathbf{w}_t^{\mathrm{in}},
\]
where \(\alpha\in[0,1]\) controls the refinement strength. Further implementation details of the inner architecture and refinement construction are provided in the supplementary material.

The inner reward measures the incremental contribution of daily refinement relative to the active base portfolio:
\[
r_t^{\mathrm{in}}
=
\log\!\left(\mu_t\,\mathbf{y}_t^\top \mathbf{w}_t\right)
-
\log\!\left(\mathbf{y}_t^\top \mathbf{b}_t\right),
\]
up to the constant reward scale used by PPO. The first term is the actual one-step executed log return after transaction cost, and the second term is the one-step log return of the active base portfolio. By measuring refinement value relative to the current base portfolio rather than in absolute terms, this reward encourages the inner actor to extract short-horizon incremental gains beyond the existing base portfolio while remaining cost-aware.


\subsection{Controller-Guided Strategic Revision}

Beyond deciding \emph{what} base portfolio to hold and \emph{how} to perform daily refinement, \textit{CMTFlow} must also determine \emph{when} the active base portfolio should be revised. Fixed holding periods are inadequate: delayed revision may leave the portfolio tied to an obsolete base, whereas overly frequent revision can undermine base-portfolio stability and increase transaction cost. The controller therefore acts as a trainable event gate: it does not directly allocate capital, but decides whether to keep the current base portfolio or switch to the candidate base portfolio proposed by the outer actor.

\paragraph{Controller State and Exit Probability.}
At each day \(t\), the controller observes both the current holding state and the candidate replacement state. Let \(\tilde{\mathbf{w}}_t\) denote the drifted executed portfolio and \(\mathbf{w}_t^{\mathrm{cand}}\) denote the candidate base portfolio generated by the outer actor. The controller normalizes them as \(\bar{\mathbf{w}}_t=\mathrm{norm}(\tilde{\mathbf{w}}_t)\) and \(\bar{\mathbf{w}}_t^{\mathrm{cand}}=\mathrm{norm}(\mathbf{w}_t^{\mathrm{cand}})\). Its state is
\[
s_t^{\mathrm{ctrl}}
=
\{\mathbf{X}_t^{\mathrm{ctrl}},\bar{\mathbf{w}}_t,\bar{\mathbf{w}}_t^{\mathrm{cand}},\mathbf{u}_t,\mathbf{a}_t^{\mathrm{ctrl}}\}.
\]
It uses the last \(H_{\mathrm{ctrl}}\) days of a causal market window \(\mathbf{X}_{t}^{\mathrm{ctrl}}\), a compact holding-state vector
\[
\mathbf{u}_t
=
\begin{bmatrix}
d_t/H_{\mathrm{ref}},\;1-d_t/H_{\mathrm{ref}},\;
\tanh(R_t^{\mathrm{seg}}/s_R),\\
\tanh(D_t^{\mathrm{seg}}/s_D),\;
\|\bar{\mathbf{w}}_t\|_2^2
\end{bmatrix},
\]
and candidate action features
\[
\begin{aligned}
\mathbf{a}^{\mathrm{ctrl}}_t
&=
\bigl[
\Delta_t^{\mathrm{turn}},\;
\Delta_t^{\mathrm{conc}},\;
\Delta_t^{\mathrm{ovlp}}
\bigr],\\
\Delta_t^{\mathrm{turn}}
&=\|\bar{\mathbf{w}}_t^{\mathrm{cand}}-\bar{\mathbf{w}}_t\|_1,\\
\Delta_t^{\mathrm{conc}}
&=\|\bar{\mathbf{w}}_t^{\mathrm{cand}}\|_2^2,\\
\Delta_t^{\mathrm{ovlp}}
&=\sum_i \min(\bar{w}_{i,t}^{\mathrm{cand}},\bar{w}_{i,t}),
\end{aligned}
\]
which summarize candidate turnover, candidate concentration, and support overlap with the current holding. Here \(d_t\) is the age of the current holding segment, \(H_{\mathrm{ref}}\) is the fixed reference horizon used to normalize holding age, \(R_t^{\mathrm{seg}}\) is the segment return, \(D_t^{\mathrm{seg}}\) is the segment drawdown, and \(s_R,s_D\) are scaling constants.

The candidate state is encoded explicitly rather than treated as an external post-processing option. The controller first applies a one-layer asset-wise LSTM to the recent market window. It then uses two temporal attention blocks: a portfolio-state query extracts context from the whole asset sequence, and a second per-asset attention refines this context before current-holding and candidate summaries are formed as
\[
\mathbf{m}^{\mathrm{hold}}_t=\sum_i \bar{w}_{i,t}\mathbf{e}_{i,t},
\quad
\mathbf{m}^{\mathrm{cand}}_t=\sum_i \bar{w}_{i,t}^{\mathrm{cand}}\mathbf{e}_{i,t},
\]
with analogous attention summaries \(\mathbf{r}^{\mathrm{hold}}_t\) and \(\mathbf{r}^{\mathrm{cand}}_t\). The final controller feature is
\[
\begin{aligned}
\Delta\mathbf{m}_t&=\mathbf{m}^{\mathrm{cand}}_t-\mathbf{m}^{\mathrm{hold}}_t,\\
\Delta\mathbf{r}_t&=\mathbf{r}^{\mathrm{cand}}_t-\mathbf{r}^{\mathrm{hold}}_t,\\
\mathbf{v}^{\mathrm{ctrl}}_t
&=
\bigl[
\mathbf{m}^{\mathrm{hold}}_t,\;
\mathbf{r}^{\mathrm{hold}}_t,\;
\Delta\mathbf{m}_t,\;
\Delta\mathbf{r}_t,\;
\mathbf{u}_t,\;
\mathbf{a}^{\mathrm{ctrl}}_t
\bigr].
\end{aligned}
\]
Thus, the controller does not only judge whether the current holding has deteriorated; it also compares that holding with the specific candidate portfolio proposed by the outer actor. A multilayer head maps \(\mathbf{v}^{\mathrm{ctrl}}_t\) to a base exit logit \(e_t^{\mathrm{base}}\) and an auxiliary switch-advantage estimate \(\widehat{A}^{\mathrm{sw}}_t\). In the final configuration, the switch-advantage estimate is used as a bounded logit modulation,
\[
e_t
=
e_t^{\mathrm{base}}
+\eta\tanh(\widehat{A}^{\mathrm{sw}}_t/c_A),
\]
where \(\eta\) and \(c_A\) control the strength and scale of switch-advantage logit modulation.
The switch probability is
\[
\pi_t^{\mathrm{exit}}=\sigma(e_t).
\]
During stochastic training, the hold/switch action \(g_t\in\{0,1\}\) is sampled from a categorical distribution induced by this logit; during evaluation, the controller switches when \(\pi_t^{\mathrm{exit}}\) exceeds the evaluation threshold under the same daily decision protocol used in the final joint finetuning stage.

\paragraph{Controller Decisions and Segment Boundaries.}
In the final evaluation protocol, the trained controller is checked daily. Let \(d_t=t-\tau_k\) be the age of the current holding segment. Unlike the fixed actor-training schedule, the evaluation controller does not impose a minimum holding lock: after the initial position is formed, it can make a free hold/switch decision at every trading day. The reported run keeps the fixed reference horizon \(H_{\mathrm{ref}}=30\) as a maximum holding cap, so a segment is revised either when the learned exit probability crosses the threshold or when the cap is reached:
\[
\tau_{k+1}\leftarrow t
\quad\text{if}\quad
\pi_t^{\mathrm{exit}}\ge 0.5
\;\text{or}\;
d_t\ge H_{\mathrm{ref}}.
\]
This design leaves early revision timing to the learned controller rather than to a rule-based calendar schedule. In the final experiments, the controller is evaluated with a decision stride of one trading day, a probability threshold of 0.5, no minimum-hold constraint, and a 30-day maximum-hold cap. Fixed 5/10/20/30/60-day controllers are used only as ablation baselines.

\paragraph{Controller Training Objective.}
The outer and inner actors are first warmed up under fixed holding segments so that they provide a stable HRL backbone. The controller is then trained with counterfactual policy gradients over controller-active rollouts, without external switch labels. For a rollout \(\xi\), let \(C\) denote the controlled path induced by sampled controller decisions and \(B\) denote the fixed-segment reference path. The controller reward is defined from relative log-return improvement with a penalty for excessive switching:
\[
R_{\xi}^{\mathrm{ctrl}}
=
\lambda_R\bigl(\log V_{\mathrm{end}}^{C}-\log V_{\mathrm{end}}^{B}\bigr)
-
\lambda_N\,\Omega(N_C),
\]
where \(V_{\mathrm{end}}^{C}\) and \(V_{\mathrm{end}}^{B}\) are terminal values over the controller rollout, \(N_C\) is the controlled segment/switch count in the rollout, and \(\Omega(\cdot)\) penalizes overflow beyond the allowed switching budget. The stochastic controller is optimized with a REINFORCE-style objective
\[
\mathcal{L}_{\mathrm{ctrl}}^{\mathrm{pg}}
=
-\mathbb{E}_{\xi}\!\left[
R_{\xi}^{\mathrm{ctrl}}
\sum_{t\in\xi}\log\pi_{\theta^{\mathrm{ctrl}}}(g_t\mid s_t^{\mathrm{ctrl}})
\right].
\]
Auxiliary heads predict remaining-horizon holding return, remaining-horizon holding risk, and switch advantage from the same controller feature. Their losses are trained from rollout statistics and counterfactual switch targets,
\[
\mathcal{L}_{\mathrm{ctrl}}
=
\mathcal{L}_{\mathrm{ctrl}}^{\mathrm{pg}}
+\lambda_r\mathcal{L}_{r}
+\lambda_d\mathcal{L}_{d}
+\lambda_a\mathcal{L}_{\mathrm{sw}},
\]
where \(\lambda_r\), \(\lambda_d\), and \(\lambda_a\) weight the return, risk, and switch-advantage auxiliary losses, respectively, and \(\mathcal{L}_{\mathrm{sw}}\) is derived from counterfactual switch targets rather than from external switch labels. In the final configuration, this switch-advantage estimate also provides the bounded logit modulation shown above, while the return and risk heads remain auxiliary representation losses. After controller policy-gradient training, a short controller-active joint finetune updates the controller, outer actor, and inner actor together with a small learning-rate multiplier, matching the final evaluation protocol.

This controller links the three layers of \textit{CMTFlow}: it preserves the outer actor's base portfolio when the current holding state remains useful, invokes a new outer decision when the state deteriorates or the candidate replacement is sufficiently attractive, and gives the inner actor a persistent base around which daily relative-return refinements can be made.

The complete training recipe has five stages. We first warm up the outer actor on fixed 30-day reference segments, warm up the inner actor for daily support-constrained refinement, and stabilize the two actors as a fixed-segment HRL backbone. We then optimize the controller by policy gradient over controller-active counterfactual rollouts without external switch labels. Finally, the best controller checkpoint initializes a short controller-active end-to-end finetune in which the controller, outer actor, and inner actor are all trainable under the same daily decision protocol used at test time. The supplementary material gives pseudocode and implementation details.

\section{Experiments}
\subsection{Datasets and Experimental Setup}

We evaluate \textit{CMTFlow} on two equity markets using Nasdaq-100 and CSI-300 constituent universes. All experiments use chronological splits and retain only stocks with continuous records in the selected universe, resulting in 39 Nasdaq-100 stocks and 53 CSI-300 stocks. The Nasdaq-100 split is 2000/04/07--2017/12/29 for training, 2018/01/02--2020/04/22 for validation, and 2020/04/23--2025/10/03 for testing. The CSI-300 split is 2000/04/07--2017/12/28 for training, 2018/01/02--2019/12/31 for validation, and 2020/01/02--2025/02/28 for testing.

All methods are evaluated under a close-to-close, long-only, fully invested protocol with proportional transaction costs. Market features are normalized causally using only past observations. For \textit{CMTFlow}, both markets use transaction cost \(5\times10^{-5}\), outer/inner/controller windows of 60/10/30 days, Top-\(K=10\), refinement strength \(\alpha=0.6\), a fixed 30-day HRL warmup reference, daily controller evaluation with threshold 0.5, no minimum-hold lock, and a 30-day maximum-hold cap. The controller is trained by counterfactual policy gradient over controller-active rollouts, and the final checkpoint is obtained by a short controller-active joint finetune of the controller, outer actor, and inner actor.

\subsection{Baselines and Metrics}

We compare \textit{CMTFlow} with traditional online portfolio strategies and deep RL baselines under matched long-only, cost-aware evaluation records whenever daily traces are available. Traditional and online methods include Buy\&Hold, Markowitz \citep{MVO}, UCRP \citep{UCRP}, Anticor \citep{ANTICOR}, OLMAR \citep{OLMAR}, and WMAMR \citep{WMAMR}. Deep RL baselines include AlphaStock \citep{alphastock}, DeepAries \citep{deeparies2025}, and DeepTrader \citep{6}. To avoid mixing incompatible runs, Table~\ref{tab:main} reports only baseline records whose final metrics can be matched to the same market split and evaluation protocol. Cumulative-wealth plots require the stricter condition that a stored or replayed daily portfolio path reproduces the corresponding metric within tolerance; we include these plots in the supplementary material.

Let \(\{V_t\}_{t=0}^{T}\) denote the portfolio value sequence and \(r_t=V_t/V_{t-1}-1\) the daily return. We report four comparison metrics: total return (TR), annualized Sharpe ratio, maximum drawdown (MDD), and the Calmar-style return-to-drawdown ratio \(\mathrm{CR}=r_{\mathrm{ann}}/\mathrm{MDD}\). Annualized return is used only to compute CR, annualized volatility is used only to compute Sharpe, and neither is reported as a standalone comparison column.

\subsection{Main Results}

\begin{table}[H]
\centering
\caption{Main comparison with matched baselines. Values are percentages except Sharpe and CR. The best result in each market-metric column is bolded, and the second best is underlined.}
\label{tab:main}
\scriptsize
\resizebox{\columnwidth}{!}{
\begin{tabular}{llrrrrrrrr}
\toprule
\multirow{2}{*}{\textbf{Type}} & \multirow{2}{*}{\textbf{Model}}
& \multicolumn{4}{c}{\textbf{Nasdaq-100}}
& \multicolumn{4}{c}{\textbf{CSI-300}} \\
\cmidrule(lr){3-6}\cmidrule(lr){7-10}
& & TR\(\uparrow\) & Sharpe\(\uparrow\) & MDD\(\downarrow\) & CR\(\uparrow\)
& TR\(\uparrow\) & Sharpe\(\uparrow\) & MDD\(\downarrow\) & CR\(\uparrow\) \\
\midrule
Ours & \textit{CMTFlow} & \best{265.53} & \second{1.15} & \second{18.62} & \second{1.42}
& \second{204.99} & \best{1.14} & \second{22.78} & \second{1.09} \\
\midrule
\multirow{6}{*}{Traditional} & Anticor   & 259.97 & 0.84 & 44.59 & 0.67 & 123.82 & 0.61 & 43.33 & 0.55 \\
& Buy\&Hold & 160.64 & 0.95 & 29.79 & 0.66 & 54.37 & 0.53 & 29.86 & 0.36 \\
& Markowitz & 57.76 & 0.57 & 24.73 & 0.40 & 50.53 & 0.45 & 36.11 & 0.32 \\
& OLMAR     & 157.77 & 0.63 & 45.48 & 0.56 & 180.46 & 0.69 & 45.34 & 0.67 \\
& UCRP      & 170.39 & 1.02 & 26.63 & 0.76 & 72.58 & 0.66 & 23.56 & 0.55 \\
& WMAMR     & \second{264.39} & 0.85 & 33.88 & 0.88 & 117.12 & 0.58 & 50.52 & 0.49 \\
\midrule
\multirow{3}{*}{Deep RL} & AlphaStock & 185.35 & 1.02 & 21.90 & 0.98 & 128.39 & 0.87 & 36.37 & 0.53 \\
& DeepAries  & 162.96 & \best{1.50} & \best{10.83} & \best{1.78} & 148.76 & 0.90 & \best{15.23} & \best{1.38} \\
& DeepTrader & 196.27 & 0.90 & 24.54 & 0.94 & \best{212.81} & \second{1.03} & 31.86 & 0.83 \\
\bottomrule
\end{tabular}
}
\end{table}

Table~\ref{tab:main} shows that \textit{CMTFlow} improves the risk-return frontier rather than optimizing a single metric in isolation. On Nasdaq-100, \textit{CMTFlow} obtains the highest cumulative return among the matched methods, reaching 265.53\%, while reducing maximum drawdown to 18.62\%. This drawdown is much lower than high-return traditional baselines such as WMAMR and Anticor, whose maximum drawdowns are 33.88\% and 44.59\%, respectively. DeepAries achieves a higher Sharpe ratio, lower MDD, and the best CR on Nasdaq-100, but its total return is 162.96\%, substantially lower than \textit{CMTFlow}. \textit{CMTFlow} obtains the second-best CR of 1.42, indicating strong annual-return efficiency per unit of drawdown. Thus, the Nasdaq result should be interpreted as a balanced high-return and controlled-drawdown outcome, not as dominance on every single metric.

On CSI-300, \textit{CMTFlow} achieves 204.99\% total return, slightly below DeepTrader's 212.81\%, but improves Sharpe ratio from 1.03 to 1.14 and reduces maximum drawdown from 31.86\% to 22.78\%. Its CR is 1.09, second only to DeepAries' 1.38 among matched baselines. This market is therefore an important stress test for the claim: \textit{CMTFlow} is not simply the highest-terminal-wealth strategy in every setting, but it improves the quality and stability of the investment path. Overall, the main results support the claim that explicit base-revision control can move the strategy toward a more favorable return--risk trade-off.

\subsection{Ablation Study}

The ablation study identifies the controller as the main source of adaptive behavior. On Nasdaq-100, adding the controller to the outer actor improves total return from 220.42\% to 237.50\%, reduces MDD from 32.09\% to 21.24\%, and raises CR from 0.74 to 1.18; the full model further reaches 265.53\% return, 18.62\% MDD, and the best CR of 1.42. On CSI-300, Outer + Controller achieves 237.77\% return, 1.22 Sharpe, and 1.16 CR, while the full model keeps a slightly lower MDD and outperforms all fixed-window controllers in total return. This contrast clarifies the roles of the modules: the controller is the dominant mechanism for dynamic revision, whereas the inner actor is a local weight-refinement module whose return contribution is market-dependent rather than uniformly positive.

The fixed-window results show that the gain is not simply caused by selecting a different constant rebalancing frequency. Fixed 5/10/20/30/60-day controllers fail to reproduce the full model's return on either market, and their lower-MDD cases on CSI-300 come with much lower terminal return. The learned controller therefore contributes state-dependent revision timing rather than a manually chosen holding period.

\subsection{Interpretability of Controller and Inner Actor}

We use interpretability analyses to test whether the two adaptive modules behave consistently with their assigned roles. For the controller, the strongest evidence is not a single switch event, but the combination of fixed-window comparisons, switch-level counterfactuals, and probability diagnostics. Compared with fixed 5/10/20/30/60-day holding windows, the learned controller achieves the best TR, Sharpe, and CR on both markets among the fixed-window group, supporting the claim that it learns state-dependent revision timing rather than a constant holding period.

Switch-level counterfactuals give a conservative decision-level check. For every free switch, we freeze the switch-date alternatives and compare the controller-selected base with the old-base continuation over the same remaining horizon, so later switches do not contaminate the comparison. Across all free switches, the mean switch and old-holding returns are close: 3.29\% versus 3.24\% on Nasdaq-100, and 4.43\% versus 4.39\% on CSI-300. We therefore do not claim large average per-switch alpha; rather, switching is not systematically worse than continuing, and path-level gains arise when maintenance decisions prevent exposure to important deteriorating windows. The exit probability is likewise a nonlinear policy signal, not a one-dimensional future-return predictor: its global linear correlation with subsequent switch advantage is weak (\(-0.05\) and 0.01), as expected from a decision depending jointly on holding state, candidate turnover, support overlap, costs, and market regime.

For the inner actor, the meaningful object is the relative tilt \(\Delta w_{i,t}^{\mathrm{inner}}=w_{i,t}-b_{i,t}\), not an independent stock-selection decision. In selected windows, this tilt aligns with subsequent cross-sectional relative performance: the tilt--future-return correlation is 0.46 on Nasdaq-100 and 0.33 on CSI-300, with positive alignment on 73.33\% and 70.00\% of days, respectively. This supports a modest but interpretable role: the inner actor provides within-holding-period weight flexibility, while the controller remains the dominant adaptive revision mechanism.

\FloatBarrier
\raggedbottom
\bibliography{aaai2026}

% Check whether the conference requires a reproducibility checklist to be included in the paper.
% If so, you can uncomment the following line and ajust the path to include it.
% \input{../../ReproducibilityChecklist/LaTeX/ReproducibilityChecklist.tex}

\end{document}

\appendix

\section{Additional Details of Hierarchical Policies}
\label{app:hier_policy}

\subsection{Outer Architecture and Base Portfolio Construction}
\label{app:hier_outer}

This section provides additional implementation details of the outer policy, whose role is to construct a segment-level candidate base portfolio at each controller decision point or fixed warmup boundary.

At decision time \(t=\tau_k\), the outer actor observes
\[
s^{\mathrm{out}}_{\tau_k}
=
\bigl\{\mathbf{X}^{\mathrm{out}}_{\tau_k},\tilde{\mathbf{w}}_{\tau_k}\bigr\},
\]
where \(\mathbf{X}^{\mathrm{out}}_{\tau_k}\in\mathbb{R}^{M\times L_{\mathrm{out}}\times F}\) denotes the long-horizon market observation tensor for the \(M\) candidate assets, \(L_{\mathrm{out}}\) is the outer lookback window, and \(F\) is the feature dimension. The vector \(\tilde{\mathbf{w}}_{\tau_k}\in\mathbb{R}^{M}\) is the drifted executed portfolio at the current decision point and provides the outer actor with the current portfolio context.

The outer actor first applies an asset-wise temporal encoder to \(\mathbf{X}^{\mathrm{out}}_{\tau_k}\) to extract long-horizon sequential representations:
\[
\mathbf{H}^{\mathrm{out}}_{\tau_k}
=
\mathrm{LSTM\mbox{-}HA}\!\left(\mathbf{X}^{\mathrm{out}}_{\tau_k}\right),
\]
where \(\mathbf{H}^{\mathrm{out}}_{\tau_k}\in\mathbb{R}^{M\times d_{\mathrm{out}}}\) is the resulting asset-wise latent representation and \(d_{\mathrm{out}}\) is the outer hidden dimension. The LSTM-HA module summarizes asset-specific temporal patterns over the outer horizon.

To incorporate cross-asset dependency, the latent representation is then processed by the cross-asset attention allocator (CAAN):
\[
\mathbf{A}^{\mathrm{out}}_{\tau_k}
=
\mathrm{CAAN}\!\left(
\mathbf{H}^{\mathrm{out}}_{\tau_k}
\right),
\]
where \(\mathbf{A}^{\mathrm{out}}_{\tau_k}\in\mathbb{R}^{M\times d_{\mathrm{out}}}\) contains cross-asset contextual representations. These representations are then fused with the most recent cross-sectional feature matrix \(\mathbf{x}^{\mathrm{last}}_{\tau_k}\in\mathbb{R}^{M\times F}\) and the drifted portfolio context:
\[
\mathbf{C}^{\mathrm{out}}_{\tau_k}
=
\mathrm{Fusion}^{\mathrm{out}}\!\left(
\mathbf{A}^{\mathrm{out}}_{\tau_k},
\mathbf{x}^{\mathrm{last}}_{\tau_k},
\tilde{\mathbf{w}}_{\tau_k}
\right).
\]
This fusion allows the outer actor to allocate capital according to asset-wise temporal signals, cross-asset context, and the current portfolio state.

A stochastic policy head then maps \(\mathbf{C}^{\mathrm{out}}_{\tau_k}\) to the mean and log-standard-deviation of an asset-wise Gaussian allocation distribution:
\[
\boldsymbol{\mu}^{\mathrm{out}}_{\tau_k},
\log \boldsymbol{\sigma}^{\mathrm{out}}_{\tau_k}
=
\mathrm{MLP}^{\mathrm{out}}\!\left(\mathbf{C}^{\mathrm{out}}_{\tau_k}\right).
\]
A raw allocation score vector is sampled as
\[
\hat{\mathbf{q}}^{\mathrm{out}}_{\tau_k}
\sim
\mathcal{N}\!\left(
\boldsymbol{\mu}^{\mathrm{out}}_{\tau_k},
\mathrm{diag}\bigl((\boldsymbol{\sigma}^{\mathrm{out}}_{\tau_k})^2\bigr)
\right),
\]
and then squashed by \(\tanh(\cdot)\) to obtain bounded preference scores:
\[
\mathbf{q}^{\mathrm{out}}_{\tau_k}
=
\tanh\!\left(\hat{\mathbf{q}}^{\mathrm{out}}_{\tau_k}\right).
\]

To construct a sparse long-only candidate base portfolio, only the Top-\(K\) assets with the largest scores are retained. Let \(\mathcal{I}_{\tau_k}=\mathrm{TopK}(\mathbf{q}^{\mathrm{out}}_{\tau_k},K)\) denote the index set of the selected assets. The candidate base portfolio is then given by
\[
w_{i,\tau_k}^{\mathrm{cand}}
=
\begin{cases}
\dfrac{\exp(q_{i,\tau_k}^{\mathrm{out}})}
{\sum_{j\in\mathcal{I}_{\tau_k}}
\exp(q_{j,\tau_k}^{\mathrm{out}})},
& i\in\mathcal{I}_{\tau_k},\\[10pt]
0, & \text{otherwise}.
\end{cases}
\]
Here the exponential transformation ensures positive portfolio weights, and normalization guarantees that the selected weights sum to one. Thus, the outer actor converts long-horizon market context into a sparse segment-level candidate base portfolio.

\subsection{Inner Architecture and Support-Constrained Refinement}
\label{app:hier_inner}

This section provides additional implementation details of the inner policy, whose role is to refine execution daily around the active base portfolio.

At each trading day \(t\) within the current segment, the inner actor observes
\[
s_t^{\mathrm{in}}
=
\bigl\{\mathbf{X}^{\mathrm{in}}_t,\tilde{\mathbf{w}}_t,\mathbf{b}_t\bigr\},
\]
where \(\mathbf{X}^{\mathrm{in}}_t\in\mathbb{R}^{M\times L_{\mathrm{in}}\times F}\) denotes the short-horizon market observation tensor, \(\tilde{\mathbf{w}}_t\) is the current drifted executed portfolio, and \(\mathbf{b}_t\) is the active base portfolio selected by the controller. Compared with the outer actor, the inner actor uses a much shorter lookback window \(L_{\mathrm{in}}\), because its role is to react to local fluctuations rather than to construct stable medium-horizon exposure.

The inner actor first applies an asset-wise LSTM-attention encoder to extract short-horizon representations:
\[
\mathbf{H}^{\mathrm{in}}_t
=
\mathrm{LSTM\mbox{-}Attn}\!\left(\mathbf{X}^{\mathrm{in}}_t\right),
\]
where \(\mathbf{H}^{\mathrm{in}}_t\in\mathbb{R}^{M\times d_{\mathrm{in}}}\) is the inner latent representation and \(d_{\mathrm{in}}\) is the inner hidden dimension. The implementation uses a two-layer LSTM for each asset window and two temporal attention layers that use the latest hidden state as the query, allowing the inner actor to focus on the most relevant days inside the short lookback window.

A stochastic policy head then conditions on \(\mathbf{H}^{\mathrm{in}}_t\), the current drifted executed portfolio, and the active base portfolio to parameterize an asset-wise Gaussian refinement distribution:
\[
\boldsymbol{\mu}^{\mathrm{in}}_t,\log \boldsymbol{\sigma}^{\mathrm{in}}_t
=
\mathrm{MLP}^{\mathrm{in}}\!\left(
\mathbf{H}^{\mathrm{in}}_t,
\tilde{\mathbf{w}}_t,
\mathbf{b}_t
\right).
\]
A raw refinement signal is sampled as
\[
\hat{\mathbf{q}}^{\mathrm{in}}_t
\sim
\mathcal{N}\!\left(
\boldsymbol{\mu}^{\mathrm{in}}_t,
\mathrm{diag}\bigl((\boldsymbol{\sigma}^{\mathrm{in}}_t)^2\bigr)
\right),
\]
and used directly as the masked softmax logit:
\[
\mathbf{q}^{\mathrm{in}}_t
=
\hat{\mathbf{q}}^{\mathrm{in}}_t.
\]

To preserve the structural role of the outer actor, the inner actor is constrained to refine only the support of the active base portfolio. Let
\[
m_{i,t}
=
\mathbb{I}\!\left(b_{i,t}>0\right)
\]
be the support mask induced by \(\mathbf{b}_t\). The masked refinement signal is
\[
\bar{q}_{i,t}^{\mathrm{in}}
=
m_{i,t}\exp(\hat{q}_{i,t}^{\mathrm{in}}),
\]
and the final refinement portfolio is obtained by normalization over the active support:
\[
w_{i,t}^{\mathrm{in}}
=
\frac{\bar{q}_{i,t}^{\mathrm{in}}}
{\sum_{j=1}^{M}\bar{q}_{j,t}^{\mathrm{in}}}.
\]
Hence, the inner actor can only redistribute portfolio mass within the support of the current base portfolio and cannot arbitrarily introduce new assets outside that support. This is precisely what keeps daily refinement anchored to the outer actor's segment-level structural decision.

The final executed portfolio is obtained by convex combination of the active base portfolio and the inner refinement:
\[
\mathbf{w}_t
=
(1-\alpha)\mathbf{b}_t
+
\alpha\mathbf{w}^{\mathrm{in}}_t,
\]
where \(\alpha\in[0,1]\) controls the strength of daily refinement.

\section{Additional Details of End-to-End Staged Training}
\label{app:hier_training}

To stabilize hierarchical optimization while preserving the final daily controller policy, \textit{CMTFlow} is trained with the same staged end-to-end schedule used in the experiment script. Algorithm~\ref{alg:cmtflow_training} gives the compact pseudocode in the main text, while this appendix expands the rollout construction and reward definitions. During the HRL warmup stages, holding segments are constructed with a fixed reference length \(H\), i.e.,
\[
\tau_k-\tau_{k-1}=H,\qquad k=1,\dots,K.
\]
This fixed-length segmentation is introduced only for early actor-training stability. It provides the outer actor with a consistent future holding horizon and avoids introducing controller-induced boundary variation before the outer/inner backbone is usable. During final joint finetuning and deployment, by contrast, segment boundaries are determined by the learned daily controller subject to the same 30-day maximum-hold cap used in the reported evaluation.

Under fixed-length HRL warmup, the outer actor is invoked once at each segment boundary \(t=\tau_k\), while the inner actor acts daily within the resulting segment. Let \(\mathcal{B}_{\mathrm{out}}\) denote the buffer of outer transitions and \(\mathcal{B}_{\mathrm{in}}\) denote the buffer of inner transitions collected from the training environment. The final pipeline contains five stages: outer warmup, inner warmup, fixed-HRL outer/inner joint training, counterfactual controller policy-gradient training, and controller-active end-to-end joint finetuning.

\paragraph{Stage I: Outer warmup.}
In the first stage, we optimize the outer actor alone in order to learn a stable segment-level candidate-base policy. At each fixed warmup boundary \(t=\tau_k\), the outer actor observes \(s_{\tau_k}^{\mathrm{out}}\), outputs a candidate base portfolio \(\mathbf{w}^{\mathrm{cand}}_{\tau_k}\), and receives the segment-level reward
\[
r_{\tau_k}^{\mathrm{out}}
=
\sum_{t=\tau_k}^{\tau_k+H-1}
\log\!\left(\mu_t\,\mathbf{y}_t^\top\mathbf{w}_t\right),
\]
where \(\mathbf{w}_t\) is the executed portfolio along the fixed reference segment and \(\mu_t\) applies the proportional transaction cost. In this stage, the inner actor is disabled, so the executed portfolio within each segment is determined by the outer candidate accepted as the fixed base and its passive drift. This isolates the segment-level allocation problem and allows the outer actor to learn a robust strategic anchor without interference from daily refinement. The outer parameters \(\theta^{\mathrm{out}}\) are optimized with PPO using transitions in \(\mathcal{B}_{\mathrm{out}}\), with a future-stock-return auxiliary prediction loss on the outer encoder.

\paragraph{Stage II: Inner warmup.}
In the second stage, the best outer-warmup checkpoint is used to generate a stable active base sequence \(\{\mathbf{b}_t\}\) throughout training episodes, and the inner actor is optimized to learn daily refinement around that base. At each trading day represented in \(\mathcal{B}_{\mathrm{in}}\), the inner actor observes \(s_t^{\mathrm{in}}\), outputs a refinement portfolio \(\mathbf{w}_t^{\mathrm{in}}\), and the final executed portfolio is
\[
\mathbf{w}_t
=
(1-\alpha)\mathbf{b}_t
+
\alpha \mathbf{w}_t^{\mathrm{in}}.
\]
The inner reward is
\[
r_t^{\mathrm{in}}
=
\log\!\left(\mu_t\,\mathbf{y}_t^\top \mathbf{w}_t\right)
-
\log\!\left(\mathbf{y}_t^\top \mathbf{b}_t\right).
\]
The implementation also uses a small next-return prediction auxiliary loss for the inner encoder, which encourages the LSTM-attention representation to preserve short-horizon relative-return information. This stage isolates the daily execution problem and avoids early co-adaptation between the two actor layers. The inner parameters \(\theta^{\mathrm{in}}\) are optimized with PPO using transitions in \(\mathcal{B}_{\mathrm{in}}\).

\paragraph{Stage III: Fixed-HRL joint training.}
After the two warmup stages, the outer and inner actors are jointly optimized under the same fixed 30-day reference segmentation, with the controller inactive. This stage lets the sparse base portfolio and the support-constrained daily refinement adapt to each other before the system is exposed to adaptive segment boundaries. The best fixed-HRL checkpoint is then used as the initialization for controller policy-gradient training.

\paragraph{Stage IV: Counterfactual controller policy gradient.}
The controller is trained to decide whether each candidate outer portfolio should replace the current holding state. For each controller rollout, we compare the controlled trajectory with a fixed-segment reference trajectory and assign the controller a relative investment reward. Let \(\theta^{\mathrm{ctrl}}\) denote the controller parameters. The controller objective can be written as
\[
\max_{\theta^{\mathrm{ctrl}}}
\;
\mathbb{E}\!\left[
\sum_{\xi\in\mathcal{B}_{\mathrm{ctrl}}}
r_{\xi}^{\mathrm{ctrl}}
\right],
\]
where \(\xi\) indexes controller rollouts. In the final setting, \(r_{\xi}^{\mathrm{ctrl}}\) is the coefficient-scaled log-return uplift over the fixed reference path minus a normalized overflow penalty when the controlled trajectory exceeds the configured segment budget. The current training path does not use external switch labels; switch decisions are learned from sampled controller actions and their realized counterfactual investment consequences. Auxiliary return, risk, and switch-advantage heads are trained from rollout statistics; only the switch-advantage estimate is used to modulate the switch logit, and none of these terms is an external switch label.

\paragraph{Stage V: Controller-active end-to-end joint finetuning.}
After controller policy-gradient training, we load the best controller checkpoint and run a short controller-active joint finetuning stage. In this stage the controller, outer actor, and inner actor are all trainable; rollouts use the daily controller decision protocol rather than a fixed-calendar rule; and a small joint learning-rate multiplier is applied to avoid destroying the fixed-HRL backbone. The final checkpoint is selected on validation performance under the same controller-active evaluation protocol used for the reported test results.

\paragraph{PPO optimization.}
The outer and inner actors are optimized using clipped proximal policy optimization (Clip-PPO) during actor training. For either layer, let \(\pi_{\theta}(\cdot)\) denote the current policy and \(\pi_{\theta_{\mathrm{old}}}(\cdot)\) the behavior policy used to collect trajectories. For a sampled transition \((s_t,a_t)\), the PPO probability ratio is
\[
\varrho_t(\theta)
=
\frac{\pi_{\theta}(a_t\mid s_t)}
{\pi_{\theta_{\mathrm{old}}}(a_t\mid s_t)}.
\]
The clipped surrogate objective is
\[
\mathcal{L}_{\mathrm{PPO}}
=
\mathbb{E}_t\!\left[
\min\!\bigl(
\varrho_t(\theta)\hat{A}_t,
\mathrm{clip}(\varrho_t(\theta),1-\epsilon,1+\epsilon)\hat{A}_t
\bigr)
\right],
\]
where \(\hat{A}_t\) is the estimated advantage and \(\epsilon\) is the clipping coefficient. For the outer actor, this objective is computed over segment-level transitions in \(\mathcal{B}_{\mathrm{out}}\); for the inner actor, it is computed over daily transitions in \(\mathcal{B}_{\mathrm{in}}\). Thus, although both actor layers are optimized with the same RL algorithm, they operate on different temporal resolutions and different reward signals.

Overall, the staged schedule first establishes a reliable segment-level allocator, then learns a daily refinement policy around that anchor, then learns controller switching from counterfactual rollout rewards, and finally updates all modules together under the controller-active daily decision rule. This progressive design improves optimization stability while keeping the final model end-to-end aligned with the evaluation logic used in our experiments.

\subsection{Final Evaluation and Matching Protocol}
\label{app:eval_protocol}

Table~\ref{tab:eval_protocol} summarizes the implementation details that are most relevant to reproducing the final evaluation and explanation experiments. Checkpoint selection is stage-specific in the actual script: fixed-HRL actor phases use the configured actor metrics, controller-PG uses validation return in this run, and the final controller-active joint checkpoint is initialized from the best controller-PG model and then selected under the same controller-active validation protocol. The controller is evaluated as a daily event policy with no minimum-hold lock and a 30-day maximum-hold cap.

\begin{table}[t]
\centering
\caption{Final evaluation and matching details.}
\label{tab:eval_protocol}
\footnotesize
\setlength{\tabcolsep}{3pt}
\begin{tabular}{p{0.38\columnwidth}p{0.52\columnwidth}}
\toprule
\textbf{Item} & \textbf{Setting} \\
\midrule
Final checkpoint & Selected by the controller-active validation protocol \\
Actor reference segment & 30 trading days \\
Eval rule & Daily check, threshold 0.5 \\
Holding constraint & No minimum lock; 30-day max-hold cap \\
Controller rollout setting & Length 600; 12 windows per PG batch \\
Switch budget & Max switches 30; overflow penalty 0.001 \\
Fixed-window baselines & 5, 10, 20, 30, and 60 trading days \\
Baseline matching & Plot only trajectories that reproduce table metrics within tolerance \\
\bottomrule
\end{tabular}
\end{table}

\section{Additional Details of Controller State Construction}
\label{app:ctrl_state}

This section summarizes the controller state used by the final end-to-end training path. The controller is not trained with external switch labels. Instead, every controller decision is conditioned on a causal comparison between the current holding and the candidate outer portfolio, and the decision is optimized from counterfactual rollout rewards.

At day \(t\), the controller receives the recent asset feature tensor
\[
\mathbf{X}^{\mathrm{ctrl}}_t
\in
\mathbb{R}^{M\times H_{\mathrm{ctrl}}\times F},
\]
where \(H_{\mathrm{ctrl}}\) is the controller lookback window. The tensor is constructed causally from historical market features and is encoded by the controller's asset-wise LSTM-attention encoder. Let \(\bar{\mathbf{w}}_t\) denote the normalized current holding after price drift and let \(\bar{\mathbf{w}}_t^{\mathrm{cand}}\) denote the normalized candidate base portfolio proposed by the outer actor. The controller forms weighted summaries for both portfolios:
\[
s_t^{\mathrm{ctrl}}
=
\{\mathbf{X}_t^{\mathrm{ctrl}},\bar{\mathbf{w}}_t,\bar{\mathbf{w}}_t^{\mathrm{cand}},\mathbf{u}_t,\mathbf{a}_t^{\mathrm{ctrl}}\}.
\]
\[
\mathbf{m}^{\mathrm{hold}}_t
=
\sum_i \bar{w}_{i,t}\mathbf{e}_{i,t},
\qquad
\mathbf{m}^{\mathrm{cand}}_t
=
\sum_i \bar{w}_{i,t}^{\mathrm{cand}}\mathbf{e}_{i,t},
\]
where \(\mathbf{e}_{i,t}\) is the encoded asset representation.

The controller also uses compact portfolio-state and action-comparison features:
\[
\mathbf{u}_t=
\begin{bmatrix}
d_t/H_{\mathrm{ref}}\\
1-d_t/H_{\mathrm{ref}}\\
\tanh(R_t^{\mathrm{seg}}/s_R)\\
\tanh(D_t^{\mathrm{seg}}/s_D)\\
\|\bar{\mathbf{w}}_t\|_2^2
\end{bmatrix},
\]
\[
\mathbf{a}^{\mathrm{ctrl}}_t=
\begin{bmatrix}
\|\bar{\mathbf{w}}_t^{\mathrm{cand}}-\bar{\mathbf{w}}_t\|_1\\
\|\bar{\mathbf{w}}_t^{\mathrm{cand}}\|_2^2\\
\sum_i \min(\bar{w}_{i,t}^{\mathrm{cand}},\bar{w}_{i,t})
\end{bmatrix}.
\]
The raw environment state also stores cumulative auxiliary fields for compatibility, but the controller head used in the reported model transforms it into the five-dimensional state vector \(\mathbf{u}_t\) above before forming the final switch feature.
Thus, the switch probability is a function of four coupled signals: the recent market sequence, the quality of the current holding, the properties of the candidate replacement, and the cost/overlap implied by switching. This construction is why the controller probability should be interpreted as a nonlinear policy signal rather than as a one-dimensional forecast of next-period return.
\end{document}
